\begin{notechapter}
\chapter{Adjoints, Projections and Least Squares}\label{ch:adjoint-operator}\label{ch:least-squares-geometry}
\chaptersubtitle{Why the transpose appears everywhere}

\begin{chapterintro}
The transpose $A^\top$ appears in almost every formula of machine learning: the
normal equations $A^\top A\hat x = A^\top y$, the backpropagated gradient
$W^\top\delta$, the covariance $\mathbf X^\top\mathbf X$. This chapter explains what
the transpose \emph{is}: the matrix of the adjoint operator, defined by the inner
product. With that notion, orthogonal projections and least squares become
geometric statements with short proofs. You will learn
\begin{itemize}[nosep]
  \item the difference between the dual operator (algebraic) and the adjoint
        operator (geometric), and how the Riesz theorem connects them;
  \item why the matrix of the adjoint is $A^\top$ for the standard dot product,
        and $M^{-1}A^\top N$ for other inner products;
  \item that a projection is orthogonal exactly when it is self-adjoint, also for
        a general inner product $\inner{x}{y}_M = x^\top My$;
  \item the fundamental relation $\ker(A^\top) = \operatorname{Im}(A)^\perp$ and the
        normal equations of least squares, with a worked example;
  \item when the least-squares solution exists and is unique.
\end{itemize}
\prerequisites{vector spaces, linear maps and matrices; inner products.}
\end{chapterintro}

\section{The transpose and the dot product}\label{adjoint-operator:sec:transpose}

For a real matrix $A\in\R^{n\times m}$, representing a linear map $\R^m\to\R^n$, the
transpose $A^\top\in\R^{m\times n}$ swaps rows and columns: $(A^\top)_{ji} = A_{ij}$.
This purely combinatorial definition hides a geometric property.

\begin{proposition}[Moving a matrix across a dot product]\label{adjoint-operator:prop:transpose}
For all $x\in\R^m$ and $y\in\R^n$, $\;(Ax)\cdot y = x\cdot(A^\top y)$.
\end{proposition}

\begin{proof}
$(Ax)\cdot y = (Ax)^\top y = x^\top A^\top y = x\cdot(A^\top y)$, using $(Ax)^\top = x^\top A^\top$.
In coordinates: both sides equal $\sum_{i,j}A_{ij}x_jy_i$.
\end{proof}

The rest of the chapter turns this property into a definition.

\begin{definition}[Inner product space]
An \emph{inner product} on a real vector space $V$ is a map
$\inner{\cdot}{\cdot}:V\times V\to\R$ that is symmetric, linear in each argument and
positive definite ($\inner xx>0$ for $x\neq0$). (Over $\C$, symmetry becomes
conjugate symmetry $\inner xy = \overline{\inner yx}$ and linearity holds in the first
argument only.)
\end{definition}

On $\R^p$, every inner product has the form $\inner xy_M = x^\top My$ for a symmetric
positive definite matrix $M$; the dot product is $M = I$. Other choices appear in
statistics: $M = \diag(1/s_1^2,\dots,1/s_p^2)$ makes all variables equally important
in a distance (\cref{ch:pca-overview}).

\section{Dual operator and adjoint operator}\label{adjoint-operator:sec:dual-adjoint}

\subsection{The dual operator: algebra only}

\begin{definition}[Dual space and dual operator]
The \emph{dual space} $V^*$ of a vector space $V$ is the space of linear functionals
$\phi:V\to\R$. For a linear map $T:V\to W$, the \emph{dual operator}
$T^*_{\text{dual}}:W^*\to V^*$ is defined by
\[
  T^*_{\text{dual}}(\phi) = \phi\circ T,\qquad\text{i.e.}\quad T^*_{\text{dual}}(\phi)(v) = \phi(T v).
\]
\end{definition}

The dual operator exists for every linear map and needs no inner product: it maps
\emph{functionals} on $W$ to functionals on $V$ by pre-composition. In finite
dimension, with bases of $V$, $W$ and the dual bases of $V^*$, $W^*$, its matrix is the
transpose $A^\top$ of the matrix $A$ of $T$ --- also over $\C$, where no complex
conjugation appears.

\subsection{The adjoint operator: geometry}

\begin{definition}[Adjoint]
Let $V$, $W$ be inner product spaces (Hilbert spaces in infinite dimension) and
$A:V\to W$ linear (and bounded). The \emph{adjoint} of $A$ is the unique linear map
$A^*:W\to V$ such that
\[
  \inner{Av}{w}_W = \inner{v}{A^*w}_V\qquad\text{for all }v\in V,\ w\in W.
\]
\end{definition}

The adjoint maps \emph{vectors} to vectors, in the reverse direction. Its matrix
depends on the inner products.

\begin{proposition}[Matrix of the adjoint]\label{adjoint-operator:prop:matrix}
Let $A\in\R^{n\times m}$, with inner products $\inner xy_{M} = x^\top My$ on $\R^m$ and
$\inner xy_N = x^\top Ny$ on $\R^n$. Then $A^* = M^{-1}A^\top N$. In particular, for the
standard dot products ($M = I$, $N = I$), $A^* = A^\top$.
\end{proposition}

\begin{proof}
We need $(Av)^\top Nw = v^\top M(A^*w)$ for all $v$, $w$, i.e.\ $A^\top N = MA^*$, hence
$A^* = M^{-1}A^\top N$. Uniqueness holds because $v^\top Mu = 0$ for all $v$ forces $u = 0$.
\end{proof}

So ``the adjoint is the transpose'' is true only for the standard inner product.
Over $\C^n$ with the standard Hermitian product, the adjoint is the conjugate
transpose $A^\dagger = \bar A^\top$.

\begin{example}[The adjoint of differentiation]
On functions on $[a,b]$ vanishing at the endpoints, with $\inner fg = \int_a^bf(x)g(x)\,dx$,
integration by parts gives
\[
  \inner{Df}{g} = \int_a^bf'g\,dx = \bigl[fg\bigr]_a^b - \int_a^bfg'\,dx = \inner{f}{-Dg},
\]
so the adjoint of $D = \frac{d}{dx}$ is $-D$. (Strictly, $D$ is unbounded on
$L^2([a,b])$ --- $\norm{\sin kx}$ stays bounded while $\norm{D\sin kx}$ grows like $k$ ---
so $-D$ is its \emph{formal} adjoint; the rigorous theory of unbounded operators makes
the domain part of the definition.) This is the continuous analogue of the
difference matrix and its transpose, and the reason why the adjoint of a
convolution is a correlation (Part~VII).
\end{example}

\subsection{The Riesz theorem connects the two}

\begin{theorem}[Riesz representation]\label{adjoint-operator:thm:riesz}
In a Hilbert space $V$ (in particular any finite-dimensional inner product space),
every continuous linear functional $\phi\in V^*$ is of the form
$\phi(v) = \inner vw$ for a unique $w\in V$. The map $R_V:w\mapsto\inner{\cdot}{w}$ is
an isomorphism $V\to V^*$ (conjugate-linear over $\C$).
\end{theorem}

In finite dimension the proof is short: if $e_1,\dots,e_p$ is an orthonormal basis,
$w = \sum_i\phi(e_i)e_i$ works, since both sides of $\phi(v) = \inner vw$ are linear in
$v$ and agree on the basis.

The adjoint is the dual operator ``transported'' by the Riesz isomorphisms. Starting
from $w\in W$, turn it into a functional, pull it back with the dual operator, and
turn the result back into a vector:
\[
  A^* = R_V^{-1}\circ A^*_{\text{dual}}\circ R_W .
\]
\begin{center}
\begin{tikzcd}[row sep=large, column sep=huge]
V \arrow[r, "A"] & W \arrow[d, "R_W"] \\
V^* \arrow[u, "R_V^{-1}"] & W^* \arrow[l, "A^*_{\text{dual}}"']
\end{tikzcd}
\qquad
\raisebox{-1.5em}{$\Longrightarrow\quad A^*:W\to V$ goes ``around the square''.}
\end{center}
Indeed, for $v\in V$: $R_V(A^*w)(v) = \inner{v}{A^*w} = \inner{Av}{w} = R_W(w)(Av) = \bigl(A^*_{\text{dual}}R_W(w)\bigr)(v)$.

\begin{keyidea}
For $\R^m$ and $\R^n$ with the dot product and the standard bases, the Riesz maps
are represented by identity matrices, so the dual operator and the adjoint have the
same matrix, $A^\top$. This is why the two notions are routinely conflated. The
distinction matters as soon as the inner product is not the standard one --- as in
PCA with a metric $M$ (\cref{ch:pca-overview}) or the natural gradient
(\cref{ch:gradient-tangent-space}) --- because the adjoint, unlike the dual
operator, then changes.
\end{keyidea}

\section{Projections}\label{adjoint-operator:sec:projections}

\begin{definition}[Projection]
A linear map $P:V\to V$ is a \emph{projection} (or projector) if $P^2 = P$: projecting
twice is projecting once. It projects onto $\operatorname{Im}(P)$ along $\ker(P)$,
and $V = \operatorname{Im}(P)\oplus\ker(P)$ via $v = Pv + (I-P)v$.
\end{definition}

A projection is \emph{orthogonal} if the two pieces are orthogonal:
$\inner{Px}{(I-P)y} = 0$ for all $x,y$. Otherwise it is \emph{oblique}.

\begin{proposition}[Orthogonal $\iff$ self-adjoint]\label{adjoint-operator:prop:orthogonal}
A projection $P$ is orthogonal if and only if $P^* = P$. For the inner product
$\inner xy_M = x^\top My$ on $\R^p$, this reads $MP = P^\top M$ (``$P$ is
$M$-symmetric''); for the dot product, $P^\top = P$.
\end{proposition}

\begin{proof}
($\Leftarrow$) If $P^* = P$, then
$\inner{Px}{(I-P)y} = \inner{x}{P(I-P)y} = \inner{x}{(P - P^2)y} = 0$.

($\Rightarrow$) Orthogonality says $\inner{Px}{y} = \inner{Px}{Py}$ for all $x,y$
(write $y = Py + (I-P)y$). The right-hand side is symmetric in $x$ and $y$, hence so
is the left: $\inner{Px}{y} = \inner{Py}{x} = \inner{x}{Py}$. Thus $P^* = P$.

With $\inner xy_M$, \cref{adjoint-operator:prop:matrix} gives $P^* = M^{-1}P^\top M$, and
$P^* = P$ becomes $MP = P^\top M$.
\end{proof}

\begin{proposition}[Projection onto a line, any metric]\label{adjoint-operator:prop:line}
The $M$-orthogonal projection onto $\operatorname{span}(a)$, $a\neq0$, is
\[
  P_a = a\,(a^\top Ma)^{-1}a^\top M,\qquad\text{i.e.}\quad P_ay = \frac{\inner ay_M}{\inner aa_M}\,a .
\]
More generally, onto the column space of a full-rank $A\in\R^{p\times q}$:
$P_A = A(A^\top MA)^{-1}A^\top M$.
\end{proposition}

\begin{proof}
Write $P_ay = \alpha a$. Orthogonality of the residual to $a$ requires
$\inner{a}{y - \alpha a}_M = 0$, i.e.\ $a^\top My = \alpha\,a^\top Ma$, which gives $\alpha$.
For $P_A$, write $P_Ay = A\beta$ and require $A^\top M(y - A\beta) = 0$ (orthogonality to
every column), giving $\beta = (A^\top MA)^{-1}A^\top My$. One checks $P_A^2 = P_A$ and
$MP_A = P_A^\top M$ directly.
\end{proof}

With $M = I$: $P_A = A(A^\top A)^{-1}A^\top$, the ``hat matrix'' of least squares.

\section{Least squares}\label{adjoint-operator:sec:least-squares}

Given $A\in\R^{m\times n}$ and $y\in\R^m$, the system $Ax = y$ usually has no solution
when $m>n$: $y$ is not in the column space $\operatorname{Im}(A)$. Least squares
looks for the $x$ that makes $Ax$ as close as possible to $y$:
\[
  \hat x\in\argmin_{x\in\R^n}\norm{y - Ax}^2 .
\]
Geometrically, $A\hat x$ is the point of the subspace $\operatorname{Im}(A)$ closest
to $y$: its orthogonal projection (\cref{adjoint-operator:fig:ls}).

\begin{figure}[ht]
\centering
\begin{tikzpicture}[scale=1.15, >=Stealth]
  \fill[bookaccent!12] (-1.2,-0.6) -- (3.6,-0.6) -- (4.6,1.2) -- (-0.2,1.2) -- cycle;
  \draw[bookaccent!60] (-1.2,-0.6) -- (3.6,-0.6) -- (4.6,1.2) -- (-0.2,1.2) -- cycle;
  \node[bookaccent] at (3.9,0.95) {$\operatorname{Im}(A)$};
  \coordinate (O) at (0,0);
  \coordinate (P) at (2.2,0.35);
  \coordinate (Y) at (2.2,2.6);
  \draw[->, very thick] (O) -- (Y) node[above] {$y$};
  \draw[->, very thick, bookpitfall] (O) -- (P) node[below right] {$A\hat x$};
  \draw[->, thick, dashed, bookdef] (P) -- (Y) node[midway, right] {$e = y - A\hat x$};
  \draw (2.2,0.6) -- (2.45,0.6) -- (2.45,0.35);
  \fill (O) circle (1.5pt) node[below left] {$0$};
\end{tikzpicture}
\caption{Least squares: $A\hat x$ is the orthogonal projection of $y$ onto the column
space of $A$; the residual $e$ is orthogonal to that space.}
\label{adjoint-operator:fig:ls}
\end{figure}

The orthogonality of the residual to $\operatorname{Im}(A)$ is turned into equations by
the adjoint.

\begin{theorem}[The fundamental orthogonality relation]\label{adjoint-operator:thm:ker}
For $A\in\R^{m\times n}$, $\ker(A^\top) = \operatorname{Im}(A)^\perp$.
\end{theorem}

\begin{proof}
($\subseteq$) Let $A^\top z = 0$. Every element of $\operatorname{Im}(A)$ is $Ax$ for some
$x$, and $\inner{Ax}{z} = \inner{x}{A^\top z} = \inner{x}{0} = 0$. So $z\perp\operatorname{Im}(A)$.

($\supseteq$) Let $w\perp\operatorname{Im}(A)$. Then for every $v\in\R^n$,
$0 = \inner{Av}{w} = \inner{v}{A^\top w}$. The only vector orthogonal to all of
$\R^n$ is $0$ (take $v = A^\top w$), so $A^\top w = 0$.
\end{proof}

\begin{theorem}[Normal equations]\label{adjoint-operator:thm:normal}
$\hat x$ minimizes $\norm{y - Ax}^2$ if and only if
\[
  A^\top A\hat x = A^\top y .
\]
Solutions always exist, and $A\hat x$ --- the orthogonal projection of $y$ onto
$\operatorname{Im}(A)$ --- is the same for all of them.
\end{theorem}

\begin{proof}
Let $p$ be the orthogonal projection of $y$ onto $\operatorname{Im}(A)$ (it exists in
finite dimension). For any $x$, $y - Ax = (y - p) + (p - Ax)$ with $y - p\perp\operatorname{Im}(A)$
and $p - Ax\in\operatorname{Im}(A)$, so by Pythagoras
$\norm{y - Ax}^2 = \norm{y-p}^2 + \norm{p - Ax}^2$, minimal exactly when $Ax = p$.
Since $p\in\operatorname{Im}(A)$, such $x$ exist. Finally, $Ax = p$ iff $y - Ax\perp\operatorname{Im}(A)$,
iff $A^\top(y - Ax) = 0$ by \cref{adjoint-operator:thm:ker}.
\end{proof}

\begin{example}[Fitting a line through three points]\label{adjoint-operator:ex:line}
Fit $y\approx\beta_0 + \beta_1t$ to the points $(t,y) = (0,1)$, $(1,2)$, $(2,2)$. Then
\[
  A = \begin{bmatrix}1&0\\1&1\\1&2\end{bmatrix},\quad
  y = \begin{bmatrix}1\\2\\2\end{bmatrix},\quad
  A^\top A = \begin{bmatrix}3&3\\3&5\end{bmatrix},\quad
  A^\top y = \begin{bmatrix}5\\6\end{bmatrix}.
\]
Solving, $\hat\beta = \frac16\begin{psmallmatrix}5&-3\\-3&3\end{psmallmatrix}\begin{psmallmatrix}5\\6\end{psmallmatrix} = \bigl(\frac76,\frac12\bigr)$.
The fitted values are $A\hat\beta = (\frac76,\frac53,\frac{13}6)$ and the residual is
$e = (-\frac16,\frac13,-\frac16)$. Check: $A^\top e = (-\frac16+\frac13-\frac16,\ \frac13-\frac13) = (0,0)$ ---
the residual is orthogonal to both columns, as the theorem requires. Its first
component says that the residuals sum to zero, a consequence of having an intercept
(\cref{measuring-performance:prop:anova}).
\end{example}

\subsection{Existence and uniqueness}

\begin{proposition}\label{adjoint-operator:prop:unique}
$\ker(A^\top A) = \ker(A)$. Hence the least-squares solution is unique iff the columns
of $A$ are linearly independent, and then $\hat x = (A^\top A)^{-1}A^\top y$ and
$A\hat x = P_Ay$ with $P_A = A(A^\top A)^{-1}A^\top$.
\end{proposition}

\begin{proof}
If $Ax = 0$ then $A^\top Ax = 0$. Conversely, if $A^\top Ax = 0$ then
$0 = x^\top A^\top Ax = \norm{Ax}^2$, so $Ax = 0$. The set of solutions of the normal
equations is $\hat x + \ker(A^\top A) = \hat x + \ker(A)$.
\end{proof}

\begin{table}[ht]
\centering
\caption{The four cases of $Ax = y$.}
\label{adjoint-operator:tab:cases}
\begin{tabular}{@{}lll@{}}
\toprule
 & $y\in\operatorname{Im}(A)$ (exact solution) & $y\notin\operatorname{Im}(A)$ (least squares) \\
\midrule
$\ker(A) = \{0\}$ & unique $x$ & unique $\hat x$ \\
$\ker(A)\neq\{0\}$ & infinitely many $x$ & infinitely many $\hat x$, one $A\hat x$ \\
\bottomrule
\end{tabular}
\end{table}

Linearly independent columns require $m\geq n$; with more unknowns than equations
($n>m$), $\ker(A)\neq\{0\}$ always. When $\ker(A)\neq\{0\}$, the projection $A\hat x$ is still
unique, and among all solutions there is one of minimum norm, $\hat x = A^+y$, given by
the Moore--Penrose pseudoinverse of \cref{ch:svd}.

\begin{pitfall}[Do not form $A^\top A$ blindly]
Mathematically, the normal equations characterize the solution. Numerically,
forming $A^\top A$ squares the condition number of $A$, so small singular values are
lost to rounding. Libraries solve least squares with a QR or SVD factorization of $A$
itself (\texttt{numpy.linalg.lstsq}).
\end{pitfall}

\begin{chaptersummary}
\begin{itemize}[nosep]
  \item The dual operator $\phi\mapsto\phi\circ A$ is algebraic; the adjoint, defined by
        $\inner{Av}{w} = \inner{v}{A^*w}$, is geometric. Riesz maps vectors to functionals and
        $A^* = R_V^{-1}A^*_{\text{dual}}R_W$.
  \item For inner products $x^\top My$ and $x^\top Ny$, $A^* = M^{-1}A^\top N$; the familiar
        $A^* = A^\top$ holds for the dot product.
  \item A projection ($P^2 = P$) is orthogonal iff it is self-adjoint; with a metric $M$,
        iff $MP = P^\top M$. The $M$-orthogonal projector on $\operatorname{Im}(A)$ is
        $A(A^\top MA)^{-1}A^\top M$.
  \item $\ker(A^\top) = \operatorname{Im}(A)^\perp$; hence least squares $\iff$ normal equations
        $A^\top A\hat x = A^\top y$. The projection $A\hat x$ always exists and is unique;
        $\hat x$ is unique iff $\ker(A) = \{0\}$.
\end{itemize}
\end{chaptersummary}

\section*{Exercises}
\addcontentsline{toc}{section}{Exercises}

\begin{exercise}\label{adjoint-operator:exr:1}
Show that $(AB)^* = B^*A^*$ and $(A^*)^* = A$ for linear maps between finite-dimensional
inner product spaces.
\solutionlink{adjoint-operator:exr:1}
\end{exercise}

\begin{exercise}\label{adjoint-operator:exr:2}
Give a $2\times2$ projection matrix that is not orthogonal, and describe geometrically
what it does.
\solutionlink{adjoint-operator:exr:2}
\end{exercise}

\begin{exercise}\label{adjoint-operator:exr:3}
In \cref{adjoint-operator:ex:line}, compute the hat matrix $P_A$ and check that
$P_Ay$ gives the fitted values, that $P_A^2 = P_A$ and that $\tr P_A = 2$.
\solutionlink{adjoint-operator:exr:3}
\end{exercise}

\begin{exercise}\label{adjoint-operator:exr:4}
Show that $\operatorname{Im}(A^\top) = \ker(A)^\perp$, and deduce that the
least-squares solution of minimum norm lies in $\operatorname{Im}(A^\top)$.
\solutionlink{adjoint-operator:exr:4}
\end{exercise}
\end{notechapter}
